\section{Conclusion}
\label{sec:conclusion}

For mechanical systems with any number of degrees of freedom, the resolution gap of Gaussian states is controlled by a single matrix formula, $\Gamma = 2\tr[(B^{1/2}CB^{1/2})^{1/2}]$. We proved that $\Gamma\sigma$ is an upper bound on the gap of the best Gaussian state, to first order in $\sigma$, and that momentum directions are free. For two coupled pendulums the formula is explicit; in thirteen systems, six of them predicted before they were computed, the best Gaussian state found without any assumed shape matches it in value to a relative $4\times10^{-8}$ and in shape. Mixtures gain less in several degrees of freedom than in one, because their advantage lives in the plane of the most stretching mode and does not extend to the other directions.

The answer to the question of the introduction is that the gap does not grow with the number of degrees of freedom as such. It grows with the number of configuration directions on which the maximal strain depends, each weighted by how strongly it depends on them. An observer focused on one unstable mode can ignore uncoupled modes entirely and all momenta at no cost, and pays for a weakly coupled mode in proportion to the square root of the coupling.

Three questions remain open. The first is Conjecture~\ref{conj:optimal}: a proof that no Gaussian state does better than the family~\eqref{eq:family}. The second is the uniform gap over all states, not only Gaussian ones, which is open already for one degree of freedom~\cite{namba2026gap}. The third concerns flows that are not of the form $|p|^2/2 + V(q)$: there the strain depends on momenta, Proposition~\ref{prop:momentum} fails, and the cost of momentum directions must enter $B$.

The same balance appears in phase-space diffusion models, where a learned score plays the role of the coarse-grained density~\cite{namba2026threshold}. A formula that says which directions of a many-dimensional state are costly and which are free is what such applications need.
